This chapter turns the evidence gathered so far into a plan for the remainder of
the doctorate. It begins by stating what has already been produced and what is
therefore no longer at risk, and derives from that the deliverables the thesis
must still provide. The planned activities, experiments, and evaluation criteria
follow, each connected to the gaps identified in the previous chapters. The
chapter then addresses the conditions under which the work must be carried out,
namely the threats to the validity of the results, the responsible use of author
prediction, the open science commitments, the project risks, and the schedule
that fits the remaining time before the defense.

\section{Current Status}

At the qualification stage, this work has already produced three
main outputs. The first is an initial empirical study on source code minimaps,
Local Binary Patterns, and traditional classifiers \cite{gebara2025iwssip}. The
second is a large-scale study with 685,906 minimaps from 100 open-source
projects and 10 programming languages \cite{gebara2025ciarp}. The third is the
MinimapSignatures tool ecosystem, which validates the feasibility of bringing
minimap analysis into IDE workflows \cite{gebara2026minimapsignatures}. To these
is added an unpublished result produced for this document: the paired
anonymization ablation of Section~\ref{sec:study3-anonymization}, which gives
the confidentiality claim of the proposal a measured basis it previously
lacked.
Table~\ref{tab:publications} summarizes this production and its editorial
status; the three completed articles are reproduced in
Appendices~\ref{apendice:iwssip} to~\ref{apendice:sbes}.

\begin{table}[ht]
\centering
\caption{Scientific production associated with this doctoral research.}
\label{tab:publications}
\begin{tabular}{p{0.40\textwidth}p{0.20\textwidth}p{0.11\textwidth}p{0.13\textwidth}}
\toprule
\textbf{Paper} & \textbf{Venue} & \textbf{Role} & \textbf{Status} \\
\midrule
Software Feature Detection Using Source Code Minimaps as Visual Signatures \cite{gebara2025iwssip} & IWSSIP 2025 (IEEE) & Main author & Published \\
Source Code Minimaps at Scale: Feature-Oriented Analysis of 100 Projects \cite{gebara2025ciarp} & CIARP 2025 (Springer LNCS 16528) & Main author & Published \\
MinimapSignatures: A Multi-IDE Ecosystem for Analysis of Visual Source Code Signatures \cite{gebara2026minimapsignatures} & SBES-Tools 2026 (CBSoft) & Main author & Accepted \\
\midrule
Explainability of visual code representations (working title) & VISAPP 2027 & Main author & Planned \\
Consolidated study on confidentiality-preserving visual code analysis (working title) & Empirical Software Engineering & Main author & Planned \\
\bottomrule
\end{tabular}
\fautor
\end{table}

The last two rows are targets, not results. The proceedings volume of CIARP
2025 appeared in 2026, which is the year recorded in the corresponding
bibliography entry. VISAPP 2027 is the venue
chosen for the explainability work developed during the research internship
abroad described in Section~\ref{sec:internship}: it is a computer vision
conference whose scope covers representation learning and texture analysis,
which matches the framing of this research, and its schedule aligns with the
internship, with a regular paper deadline expected for September 2026 according
to the current call, and the conference held in February 2027. The journal target is \emph{Empirical Software
Engineering}, chosen because the final contribution is an empirical study of a
software engineering problem rather than a purely methodological result, and
because the venue expects the kind of controlled evaluation and reproducibility
package that this work plans to deliver. The submission is planned for the
second half of 2027, once repository-disjoint evaluation and the anonymization
ablation are complete, so that the article reports consolidated evidence
instead of preliminary results.

These outputs establish a solid foundation, but the thesis is not complete. The
remaining work must deepen the methodology, broaden the task set, strengthen
evaluation, and connect the tool to developer-facing scenarios. This chapter
defines the planned research activities and the main risks.

\section{Expected Deliverables}

The final doctoral work is expected to produce a set of deliverables, not a
single isolated experiment. The main deliverables are:

\begin{enumerate}
    \item a finalized methodology for source code minimap generation,
    anonymization, cataloging, and evaluation;
    \item a curated and documented minimap dataset or dataset-generation package;
    \item benchmark results comparing handcrafted descriptors and deep visual
    models under controlled splits;
    \item an anonymization study measuring the trade-off between information
    suppression and predictive performance;
    \item an explainability study showing which visual regions support model
    decisions;
    \item an extended MinimapSignatures prototype with multi-file or
    explanation-oriented feedback;
    \item at least one evaluation of developer-facing use, even if exploratory;
    \item the final thesis text and associated reproducibility package.
\end{enumerate}

The deliverables depend on one another: the dataset and methodology feed the
models, the models feed the explanation protocols, and the explanations become
the tool feedback that the developer-centered evaluation will assess.

\section{Remaining Research Activities}

The remaining doctoral work is organized into six activities.

\subsection{A1: Refined Dataset Curation}

The large-scale dataset must be refined for additional experiments. This
includes improving label hygiene, identifying generated files, separating bot and
human contributors when possible, recording repository metadata, and creating
balanced or task-specific subsets. The catalog should support reproducible
filtering decisions.

This activity is necessary because the current dataset is large but imbalanced.
Head classes dominate some objectives, and Author labels may contain bots,
organizations, or aliases. Bots are pervasive in popular open-source projects:
about a quarter of popular GitHub repositories use them \cite{wessel2018power},
so bot-authored files are likely present in a corpus selected by popularity.
Planned curation heuristics include matching author names against known bot
naming patterns (such as the \texttt{[bot]} suffix), merging aliases by
normalized name and e-mail, and flagging accounts with anomalous commit volume
for manual inspection. The proportion of bot-attributed files will be reported,
and Author experiments will be run with and without these accounts. Better
curation will make later experiments more trustworthy.

\subsection{A2: Anonymization Evaluation}

The first installment of this activity is already complete and is reported in
Section~\ref{sec:study3-anonymization}. A paired ablation over twelve classes
and 10,923 minimaps of a large government system found no accuracy cost for
anonymization, with equivalence within two percentage points established at
$p = 2.4 \times 10^{-5}$, while the transformation itself removes 56.7\% of the
entropy per pixel. That experiment retires the uncontrolled comparison on which
the confidentiality argument previously rested, and its instruments, paired
generation, repeated cross-validation on shared splits, McNemar's test and an
equivalence test against a pre-declared margin, are reused unchanged in what
follows.

Three boundaries of that result define the remaining work. It covers a single
software system, so the ablation must be repeated on the open-source corpus,
where the dataset can also be published for external replication. It covers a
single descriptor and classifier: Local Binary Patterns may be insensitive to
exactly what anonymization removes, so the comparison must be repeated with the
convolutional backbones of the large-scale study, which learn their own filters.
And it covers only the Type objective, so Project and Author must be evaluated
as well, the latter being the objective most likely to depend on suppressed
content.

Beyond replicating the two-level comparison in these settings, Activity A2 will
extend it into a graded study over several suppression levels:

\begin{itemize}
    \item baseline (the current character-to-pixel mapping: letters and digits
    replaced by canonical values, other characters kept as raw ASCII);
    \item no anonymization (raw character codes, including identifiers and literals);
    \item comment removal (strip comment regions before rendering);
    \item punctuation normalization (map all operators and delimiters to a single
    canonical value, testing whether structure beyond punctuation still carries signal);
    \item maximum suppression (all non-whitespace characters to a single value,
    preserving only blank lines and indentation).
\end{itemize}

To avoid post hoc rationalization, the ablation design is declared in this
proposal, before execution. All suppression levels will share identical
repository-aware splits, seeds, architecture, and training budget, so that the
anonymization level is the only manipulated factor. The primary metric is
macro-F1 per objective; accuracy and per-class recall are secondary. Every
comparison against the non-anonymized rendering will be reported with both a
difference test and an equivalence test, since a non-significant difference on
its own would establish nothing. The equivalence margin remains two percentage
points, the value used and declared in the completed installment. A drop beyond
that margin will be reported as a negative result and will trigger a revision of
the confidentiality claims made in this proposal. Results for all levels will be reported regardless of outcome, and
the goal remains a performance-versus-suppression trade-off curve indicating
which level is appropriate when intellectual property protection is the
primary concern.

\subsection{A3: Imbalance-Aware Modeling}

The next modeling phase will evaluate techniques designed for imbalanced
multi-class datasets. Candidate strategies include class weighting, balanced
sampling, focal loss, logit adjustment, mixup or cutmix variants, and deferred
reweighting schedules. Evaluation will prioritize macro-F1, per-class recall,
calibration, and confusion analysis, not only accuracy.

This activity responds directly to limitations observed in the large-scale
study. The goal is to improve minority-class behavior and to avoid drawing
conclusions from head-class performance alone.

\subsection{A4: Multi-Scale and Region-Aware Minimap Analysis}

Fixed 128 by 128 minimaps are compact and convenient, but they may discard
information from long files. Future experiments should compare multiple
resolutions, patch-based inputs, pyramidal representations, and region-aware
models. Explainability results can guide this process by showing which regions
are repeatedly salient.

This activity will help answer whether minimap performance is limited by input
resolution and whether region-level analysis can support more detailed tasks,
such as identifying suspicious sections of a file rather than classifying the
whole file.

\subsection{A5: Quality and Architectural-Deviation Tasks}

The thesis should move beyond Type, Project, and Author objectives toward tasks
closer to software quality. Candidate targets include code-smell likelihood,
test-file recognition, architectural role classification, generated-file
anomaly detection, and deviation from project-specific conventions.

This activity may combine minimap signals with existing static-analysis metrics
or labels. The goal is not to claim that minimaps alone can prove quality
problems, but to evaluate whether they can provide useful weak signals for
prioritization and review.

\subsection{A6: Tool Evaluation}

MinimapSignatures has been functionally validated, but it still requires
developer-centered evaluation. Future studies may involve controlled tasks,
think-aloud sessions, interviews, or scenario-based inspections. The evaluation
should measure usefulness, understandability, trust, and perceived fit with
developer workflows.

The tool should also be extended to multi-file analysis and explanation
overlays. These features are likely to make the tool more useful than single
file-level prediction alone.

\section{Planned Experiments}

Table~\ref{tab:planned-experiments} summarizes the planned experiments, the
activity each one belongs to, and its expected outcome; the experimental
details are given in the activity descriptions above.

\begin{table}[ht]
\centering
\caption{Planned experiments for the doctoral thesis.}
\label{tab:planned-experiments}
\begin{tabular}{p{0.26\textwidth}p{0.10\textwidth}p{0.46\textwidth}}
\toprule
\textbf{Experiment} & \textbf{Activity} & \textbf{Expected output} \\
\midrule
Anonymization ablation & A2 & Trade-off curves between performance and information suppression. \\
Imbalance-aware learning & A3 & Improved macro-F1 and minority-class analysis. \\
Resolution analysis & A4 & Evidence about information loss and scalability. \\
Model comparison (additional CNN backbones, transformer variants, multitask architecture) & A4 & Stronger benchmark of architecture choices under identical repository-aware splits. \\
Quality target pilot & A5 & Evidence about feasibility for quality-related analysis. \\
Tool study with developers & A6 & Feedback on usefulness, trust, and workflow integration. \\
\bottomrule
\end{tabular}
\fautor
\end{table}

\section{Evaluation Plan}

The evaluation plan combines quantitative and qualitative evidence. Quantitative
model evaluation will use train/validation/test splits, multiple random seeds
when feasible, macro-F1, weighted-F1, accuracy, ROC-AUC, calibration measures,
and confusion matrices. When objectives involve repositories or authors, splits
must be designed to avoid leakage. For example, project-specific tasks must not
accidentally train and test on near-duplicate artifacts.

Qualitative evaluation will be used in two places. First, explainability maps
will be inspected to verify whether visual evidence is plausible and stable.
Second, tool evaluation will collect developer perceptions through interviews,
questionnaires, or observational protocols. These qualitative data will not
replace quantitative metrics; they will help interpret whether the tool's
feedback is understandable and useful.

Empirical software-engineering guidelines will inform reporting, including
explicit threats to validity and transparent description of datasets, procedures,
and limitations \cite{kitchenham2002preliminary,wohlin2012experimentation,runeson2012case}.

\section{Threats to Validity}

\subsection{Construct Validity}

Construct validity concerns whether the study measures what it claims to
measure. In this project, a major construct threat is the interpretation of
labels. Type labels may mix file extension with semantic role. Project labels
may capture boilerplate rather than meaningful project identity. Author labels
may reflect bots, generated files, or organizational accounts rather than human
style. Quality labels may be noisy if derived from static-analysis tools or
proxies.

Mitigation strategies include careful label documentation, sensitivity analysis,
manual inspection of samples, exclusion of ambiguous labels when needed, and
explicit discussion of what each objective actually represents.

\subsection{Internal Validity}

Internal validity concerns whether observed effects are caused by the method
rather than confounding factors. In minimap experiments, potential confounders
include duplicate files, generated artifacts, repository templates, file naming
patterns, and data leakage across splits. If training and test sets contain
near-identical generated files, performance may be inflated.

Mitigation strategies include content hashing, duplicate removal, repository-aware
splits, time-aware splits when possible, and ablation studies that remove
headers, comments, or generated files.

\subsection{External Validity}

External validity concerns whether results generalize beyond the evaluated
corpus. The large-scale dataset includes 100 open-source projects and 10
languages, but it may not represent proprietary systems, safety-critical
software, small scripts, educational code, or low-resource languages. Industrial
code may also have conventions not present in public repositories.

Mitigation strategies include expanding the corpus, reporting repository
characteristics, evaluating additional ecosystems, and conducting industrial or
semi-industrial case studies when access is possible.

\subsection{Conclusion Validity}

Conclusion validity concerns whether the statistical and analytical conclusions
are justified. Class imbalance, multiple comparisons, and single-run results can
lead to overconfident claims. Accuracy alone may hide poor minority-class
performance.

Mitigation strategies include reporting multiple metrics, confidence intervals
or variation across seeds when feasible, per-class results, macro-averaged
metrics, and clear distinction between preliminary and final evidence.

\section{Intellectual Property and Responsible Use}

The primary ethical motivation of this project is the protection of intellectual
property. Industrial codebases often contain strategic business information
embedded in identifiers, numeric constants, formulas, and comments. The
minimap approach is designed to enable structural analysis of source code without
exposing this information to external services or collaborators.

The character-level anonymization adopted here suppresses identifier names,
numeric literals, and string content before any minimap is transmitted or
stored externally. The paired ablation reported in
Chapter~\ref{chapter:preliminary-results} found no accuracy cost for this
suppression in the setting where it was measured; Activity A2 extends the
measurement to further corpora and models.

The Author prediction objective is a methodological probe, not a surveillance
tool. Its purpose is to demonstrate that stylistic patterns, e.g., indentation
habits, spacing conventions and operator placement, survive anonymization. This
is evidence
that the model operates on structure rather than on lexical content. Author
prediction is not proposed as a deployment target in personnel monitoring,
performance evaluation, or attribution of responsibility. Any use of this
objective outside its methodological role would be a misapplication of the
research.

For practical tool deployment, the following design principles apply:

\begin{enumerate}
    \item generate minimaps locally and transmit only the anonymized image, never
    the raw source file;
    \item make the anonymization transformation auditable and documented;
    \item allow individual objectives to be disabled at the user's discretion;
    \item communicate prediction results as probabilistic signals, not as facts;
    \item keep humans responsible for all decisions based on model output.
\end{enumerate}

The MinimapSignatures backend already implements local anonymization in the IDE
plugins: the minimap is generated client-side and only the image is sent to the
inference endpoint, not the source text. This design reflects the principle that
intellectual property protection begins at the data-collection step.

\section{Open Science Practices}

The remaining doctoral work adopts open science practices explicitly. Several
artifacts are already public: the large-scale dataset is archived on Zenodo
under a versioned DOI (Chapter~\ref{chapter:tool-support}); the training
pipelines, the backend, and the IDE plugins are available in public
repositories; and the
multitask training repository includes a model card describing the model's
inputs, outputs, intended use, and limitations.

Four practices are planned for the remaining period. First, the systematic
mapping protocol, including search strings, screening decisions, and reviewer
assignments, will be packaged as a citable artifact. Second, experimental
designs will be declared before execution, as done for the anonymization
ablation in Activity A2, so that confirmatory analyses can be distinguished
from exploratory ones. Third, each experiment will ship a replication package
with configurations, seeds, scripts, and environment specifications. Fourth,
datasets and models will be documented with structured metadata (data cards
and model cards), so that reuse does not depend on contacting the author.

These practices also serve the confidentiality argument: a claim that
anonymized minimaps preserve utility is only as strong as the ability of
independent researchers to verify it.

\section{Risks and Mitigation}

Table~\ref{tab:risks} summarizes the main research risks and mitigation actions.

\begin{table}[ht]
\centering
\caption{Research risks and mitigation strategies.}
\label{tab:risks}
\begin{tabular}{p{0.26\textwidth}p{0.33\textwidth}p{0.27\textwidth}}
\toprule
\textbf{Risk} & \textbf{Impact} & \textbf{Mitigation} \\
\midrule
Dataset imbalance dominates results & Inflated accuracy and weak minority-class performance. & Use macro metrics, reweighting, balanced subsets, and per-class analysis. \\
Anonymization removes too much signal & Models lose predictive power. & Measured at no accuracy cost on one system; Activity A2 tests whether this holds for other corpora and models. \\
Stronger suppression needed for certain contexts & Current mapping may not satisfy all intellectual-property protection requirements. & Evaluate comment removal, punctuation normalization, and maximum suppression modes in Activity A2. \\
Models learn boilerplate artifacts & Poor generalization across repositories. & Use XAI, ablations, and repository-disjoint evaluation. \\
Tool feedback is not useful to developers & Weak practical contribution. & Conduct user-centered evaluation and refine interaction design. \\
Quality labels are noisy & Invalid quality-related conclusions. & Use multiple label sources and manual validation samples. \\
\bottomrule
\end{tabular}
\fautor
\end{table}

\section{Research Internship Abroad}
\label{sec:internship}

Part of the remaining work will be carried out abroad. The internship is funded
under the CAPES doctoral sandwich programme (PDSE) and will be hosted by the
Escuela T\'ecnica Superior de Ingenier\'ia Inform\'atica of the Universidad Rey
Juan Carlos, in Madrid, Spain, under the supervision of Prof. Dr. \'Angel
S\'anchez Calle, from October 2026 to February 2027; the award letter is
reproduced in Annex~\ref{anexo:pdse}. The project registered for
the internship, \emph{Explainable Artificial Intelligence applied to Code
Visualization}, is a direct continuation of the work reported in this document.

The host group is GAVAB, the high-performance research group on algorithms
applied to computer vision and biometrics of the Universidad Rey Juan Carlos.
GAVAB's central line of work is the automatic analysis of handwritten document
images, which poses the same difficulty addressed here: extracting
discriminative structure from images whose content is thin, spatially organized
strokes rather than objects with semantic boundaries.
The texture descriptors, attribution methods, and evaluation practices developed
in that setting transfer directly to minimaps.

The internship concentrates the activities for which this expertise is most
useful. Three of the planned
activities are scheduled for this period. The dataset is to be extended beyond
the currently selected open-source projects towards a broader and more
heterogeneous corpus, covering more languages, domains, and development
practices, and enriched with additional metadata such as architectural roles and
project evolution information, which corresponds to activity A1. The evaluation
of learning models is to be broadened to vision transformers and hybrid
architectures, together with the imbalance-aware strategies of activity A3, such
as focal loss, augmentation, and reweighting, which addresses activity A4.
Finally, the attribution protocols are to be consolidated, adapting
gradient-based methods to highlight the visual cues that support model
decisions.

Two practical consequences follow for the schedule. The paper submitted to
VISAPP 2027 is prepared immediately before departure and presented during the
internship, since the conference takes place in February 2027, within the
internship period. The writing of the methodology and results chapters is
placed after the return, so that the internship period is spent on
experimental work that depends on the host group rather than on text that can be
written anywhere.

\section{Tentative Timeline}

Table~\ref{tab:timeline} presents a tentative schedule after the qualification
deposit. The schedule should be refined with the advisor and co-advisors after
the qualification exam.

\begin{table}[ht]
\centering
\caption{Tentative doctoral research timeline after qualification.}
\label{tab:timeline}
\begin{tabular}{p{0.22\textwidth}p{0.62\textwidth}}
\toprule
\textbf{Period} & \textbf{Main activities} \\
\midrule
July--August 2026 & Replicate the anonymization ablation (A2) on the open-source corpus; refine the dataset catalog (A1). \\
September 2026 & Qualification exam on September 18; incorporate the committee feedback and freeze the revised research questions and thesis structure; prepare and submit the VISAPP 2027 paper; organize the internship. \\
October 2026 -- February 2027 & \textbf{Research internship abroad (PDSE), Universidad Rey Juan Carlos.} Dataset extension (A1); imbalance-aware experiments (A3); additional architectures, including vision transformers (A4); consolidation of the explanation protocols. Presentation at VISAPP 2027 in February. \\
March--May 2027 & Repository-disjoint evaluation; consolidate held-out test results and per-class metrics; begin writing the methodology and results chapters. \\
June--August 2027 & Quality or architectural-deviation pilot (A5); extend MinimapSignatures with multi-file analysis and explanation overlays; write and submit the journal article to \emph{Empirical Software Engineering}. \\
September--October 2027 & Developer-centered evaluation of the tool (A6); continue thesis writing. \\
November--December 2027 & Consolidate results, finish the thesis text, prepare artifacts and defense material. \\
\bottomrule
\end{tabular}
\fautor
\end{table}

The schedule is compact because the core representation, the datasets, and the
preliminary studies already exist. The remaining effort concentrates on
controlled evaluation, tool extension, and writing, and it fits the eighteen
months available until the expected conclusion at the end of 2027, within the
program deadline for the doctoral defense. The five months abroad shape the
schedule without interrupting it: the activities
allocated to that period are experimental and benefit from the host group, while
the writing-intensive phases are placed before departure and after return. The
replication of the anonymization ablation was moved ahead of the internship
precisely so that its results are available when the extended architectures are
evaluated abroad, where the same comparison is repeated with convolutional and
transformer backbones.
